\section{Introduction}\label{sec:introduction}

A reconstruction of an observed history and a prediction of its future
solve different problems. Kernel Affine Hull Machines (KAHMs) construct
nonlinear reconstruction maps from local reference observations
\citep{Kumar2024KAHM}. A small reconstruction error does not establish
that the resulting representation is useful for forecasting. For example,
a reconstruction residual that vanishes along an affine line loses
position along that line, although the future target may vary with
position. Conversely, a deterministic feature cannot add information to
the history from which it is computed, but it can make useful functions
available to a restricted decoder.

This distinction raises three questions. Does the reconstruction supply
functions that a simpler matched representation cannot express? Can the
magnitude of a fitted forecast correction be controlled? Does the
resulting procedure improve prediction at the available sample size and
computational budget? Each question requires its own evidence. A
comparison of function spans concerns algebraic expressiveness; a
correction bound concerns the size of one prediction component; a
forecasting comparison concerns the complete fitted and selected procedure.

We investigate these questions through \emph{reconstruction departures}.
For each local model, a departure is the difference between its nonlinear
reconstruction and its affine projection, multiplied by a locality weight.
The decoder also receives the locality weights and weighted principal
coordinates. These common features retain principal-coordinate position
within each local model and are shared
by all reconstruction comparisons. Replacing the nonlinear map with its
affine projection makes the raw departure zero. A positive-kernel RBF
reconstruction, using the same references, kernel, and locality weights,
provides a closely matched nonlinear alternative. A residual decoder
then adds a calibrated correction to a fixed raw-history ridge forecast.

The algebraic analysis identifies both redundancy and distinctness.
Normalized transfer of KAHM reconstruction weights to reference targets
is exactly a ratio of two kernel-ridge predictions. Under a
constant-row-sum and invariant-subspace condition, KAHM and RBF departures
have the same augmented raw feature span. An irregular three-reference
construction shows that the spans can also differ. These results concern
untransformed functions. They do not equate fitted ridge predictions under
different parameterizations, nor do they establish an advantage for every
reference bank or target.

The second contribution is a uniform magnitude bound for the residual
correction. Each feature block is standardized on the fitting set and
mapped into a bounded ball before ridge regression is fitted with an
unpenalized intercept. The coefficient and correction bounds account for
the empirical residual mean and the scaling of the penalty. They apply
to bounded features generally and therefore support the specified KAHM
construction without implying a property unique to KAHM. In particular,
they do not bound the raw baseline or the total forecast error.

The principal experiment compares ten procedures over four simulation
laws, two added-noise levels, and 24 independent repetitions. Fitting,
hyperparameter selection, correction calibration, and final evaluation
use separate episode sets. The protocol and prepared inputs were recorded
locally before the final forecasts were computed; this was not a public
preregistration. Five simultaneous comparisons test a 3\% nonlinear-loss
improvement margin, and a sixth tests noninferiority to ridge on a linear
control with a 5\% tolerance.

The bounded KAHM departure improves on common local-coordinate features,
but the full prespecified improvement criterion is not satisfied.
The gain over matched RBF departures is below the required margin, and
the margin over raw RBF forecasting remains unresolved. Boosted trees
and the validation-selected non-KAHM procedure have lower aggregate
loss. Supporting controlled, natural-data, and adaptive studies show why
this distinction matters beyond the principal experiment. Their findings
are summarized in Section~\ref{sec:supporting-evidence}; their complete
methods and results appear in the companion \emph{Supplementary Studies}.

The paper thus contributes a characterization of the reconstruction
feature space, a correction bound for a specified residual decoder,
and an empirical evaluation that separates a useful reconstruction
feature from an established advantage specific to KAHM. Full proofs of
the central results and the complete principal experimental protocol
are included in Appendices~\ref{app:reconstruction-mechanism}
and~\ref{app:reconstruction-results}.
